두 방정식 $P(x)=0,\, Q(x)=0$ 의 서로 다른 실근의 개수는 
7 개, 9개이고, 집합
\[ 
A= 
\{ (x,y) \,|\, 
P(x)Q(y)=0 \mbox{ 이고 } Q(x)P(y) =0, \;
x\mbox{와 } y\mbox{는 실수} \}
\]
는 무한집합이다. 집합 $A$ 의 부분집합 
\[ B = \{ (x,y) \,|\, (x,y) \in A \mbox{ 이고 } x=y \} \]
의 원소의 개수를 $n(B)$ 라고 하면, 
이것은 $P(x),\, Q(x)$ 에 따라 변한다. 
$n(B)$ 의 최댓값을 구하라.\\


[풀이] 
$P(x)Q(y) =0$ 는 $P(x)=0 \vee Q(y)=0$ 로
$Q(x)P(y) =0$ 은 $Q(x)=0 \vee P(y)=0$ 로 바꾸어 보자.
집합 $A$ 의 조건은
\[
\begin{aligned}
& (P(x)=0 \vee Q(y)=0 ) \wedge (Q(x)=0 \vee P(y)=0 )
\\[6pt]
\Longleftrightarrow\quad
& \{ (P(x)=0 \vee Q(y)=0 ) \wedge Q(x)=0 \}
  \vee \{ (P(x)=0 \vee Q(y)=0 ) \wedge P(y)=0 \}
\\[6pt]
\Longleftrightarrow\quad
& \{ (P(x)=0 \wedge Q(x)=0 ) \vee (Q(y)=0 \wedge Q(x)=0 ) \}
  \vee \{ (P(x)=0 \wedge P(y)=0 ) \vee (Q(y)=0 \wedge P(y)=0 ) \}
\end{aligned}
\]